%======================================================================
\section{The exact statement: Lagrangian surgery and collar perturbations}\label{sec:exact}
%======================================================================
Theorem~\ref{thm:cone}(b) is an inequality-type statement (as in \cite{DLM}), and no statement of that shape can
be exact: for the figure-eight knot and $n=1$ the cone is $\CI(S^3_0(4_1),w)$, whose two generators have
$\cs$ differing by $\tfrac12$ (modulo the $\tfrac14$ ambiguity of the $w$-line), while the two generators of
$\CI(\pm\Sigma(2,3,7))$ differ by $\tfrac47$.  The exact relation has generator-dependent corrections, and they
are areas.

\subsection{The zigzag}
Let $Z_n\subset T^2$ be the Lagrangian surgery of $L_{1/(n-1)}$ and $L^w_0$ at the corner $p_n$: the two lines with
the crossing resolved, smoothed in an $\epsilon$-neighbourhood of $p_n$.  It is a closed curve in the class of
$L_{1/n}$.

\begin{lemma}[zero flux]\label{lem:flux}
$Z_n$ is Hamiltonian isotopic in $T^2$ to the line $L_{1/n}=\{a+nb=0\}$ through the origin (not to $L^w_{1/n}$).
\end{lemma}
\begin{proof}
In $\R^2$ one circuit of $Z_n$ lifts to the path $\tilde p_n\to\tilde p_n+v_1\to\tilde p_n+v_1+v_2$ with
$v_1=(n-1,-1)$, $v_2=(1,0)$, and $\tilde p_n=(-\tfrac{n-1}2,\tfrac12)$.  Two isotopic closed curves on $T^2$ are
Hamiltonian isotopic iff the signed area between them vanishes.  Against the line $a+nb=c$ the signed area is a
linear function of $c$ vanishing exactly when the line passes through the midpoints $\tilde p_n+\tfrac12v_1=(0,0)$
and $\tilde p_n+v_1+\tfrac12v_2$ of the two legs, i.e.\ for $c=0$.
\end{proof}
This is the pillowcase form of the fact that $Y_{1/n}$ carries the trivial bundle.

\begin{proposition}[cone for the zigzag, with its wrong-way block]\label{prop:zigzag}
As $\epsilon\to0$, $CF(\tilde\Gam,Z_n)$ with its action filtration has the same generators and actions as
\[
\Cone\big(CF(\tilde\Gam,L_{1/(n-1)})\xrightarrow{\ \mu^2(-,p_n)\ }CF(\tilde\Gam,L^w_0)\big),
\]
and its differential is the cone differential plus a block $\beta$ from the $L^w_0$-generators to the
$L_{1/(n-1)}$-generators, counting bigons whose boundary passes through the neck along the \emph{reflex} arc (a
corner of angle $>\pi$ at $p_n$ in the limit).  It is the cone exactly when $\beta=0$.
the first summand carrying the line action $A(x)=f_\Gamma(\tilde x)-(n-1)\tilde b^2$ of Proposition~\ref{prop:action},
the second the action $f_\Gamma(\tilde x)+2\tilde a\tilde b+\tfrac n4$ (the Kirk--Klassen primitive on $b\in\tfrac12+\Z$,
offset by $\tfrac n4$).
\end{proposition}
\begin{proof}
Generators of $CF(\tilde\Gam,Z_n)$ are those of the two summands.  Bigons between points of the same summand are
bigons for that line.  A bigon between the two summands passes through the surgery region along one of the two
resolved arcs; through the convex arc it is, for $\epsilon\to0$, a triangle with a convex corner at $p_n$, which is
$\mu^2(-,p_n)$; through the reflex arc it is an immersed bigon with a reflex corner at $p_n$, which is not a Floer
triangle and is the block $\beta$.  Neither area nor local geometry excludes $\beta$ in general.  The primitive of
$\lambda$ on $Z_n$ restricts to the two line primitives, offset by a constant, up to $O(\epsilon)$.  The constant is fixed by requiring both to be primitives of $\lambda$
on $Z_n$ agreeing at the points where $Z_n$ meets $L_{1/n}$: these are the midpoints $m_1=\tilde p_n+\tfrac12v_1$,
$m_2=\tilde p_n+v_1+\tfrac12v_2$ of the two legs (Lemma~\ref{lem:flux}), not the vertex $\tilde p_n$, and the
mismatch of the two line primitives at $m_2$ (where $b\in\tfrac12+\Z$) is $n\,b^2\equiv\tfrac n4$.  The vertex offset
$(n-1)/4$ of Lemma~\ref{lem:corner} is the wrong normalisation for the cone; the numerics confirm $\tfrac n4$ on the
$w$-leg and $0$ on the other (\S\ref{sec:tests}).
\end{proof}

\subsection{From the zigzag to the straight line}
Let $\phi_t$ be an equivariant Hamiltonian isotopy of $T^2$ with $\phi_0(Z_n)=Z_n$ and $\phi_1(Z_n)=L_{1/n}$
(Lemma~\ref{lem:flux}).  Say $\Gam$ is in \emph{minimal position} for the pair if $\phi_t(Z_n)\cap\tilde\Gam$ has
constant cardinality; by the class computation this holds iff
$|\Gam\cap L_{1/n}|=|\Gam\cap L_{1/(n-1)}|+|\Gam\cap L^w_0|$, which is the case for the figure-eight and all
$n$ (counts $2|n|=2|n-1|+2$).

\begin{theorem}[exact filtered cone, pillowcase side]\label{thm:exact}
Assume \ref{H:alex}--\ref{H:Y} and minimal position.  Then
\begin{enumerate}[label=(\alph*)]
\item the isotopy induces a bijection $x\mapsto x'$ from the generators of the cone to $\Gam\cap L_{1/n}$, and
an isomorphism of complexes $CF(\tilde\Gam,Z_n)\cong CF(\tilde\Gam,L_{1/n})$ which is the identity on generators
up to handle-slides (bigon counts change only at birth--death and handle-slide moments, and the latter act by
unitriangular changes of basis); if moreover $\beta=0$ in Proposition~\ref{prop:zigzag}, then
$\Cone(\mu^2(-,p_n))\cong CF(\tilde\Gam,L_{1/n})$;
\item the action of $x'$ is $A_{\rm cone}(x)+\delta_n(x)$, where $A_{\rm cone}$ is the cone action of
Proposition~\ref{prop:zigzag} and $\delta_n(x)=\int_{R_n(x)}\omega$ is the signed area of the region in $\R^2$
bounded by the arc of $\tilde\Gam$ from $\tilde x$ to $\tilde x'$, the segment of $Z_n$ from $\tilde x$ to the
midpoint $m$ of the leg containing $\tilde x$, and the segment of $L_{1/n}$ from $m$ to $\tilde x'$;
\item $\cs$ on the generators of $\CI(Y_{1/n})$ is $A(x')+c_i$, exactly (Proposition~\ref{prop:action}).
\end{enumerate}
\end{theorem}
\begin{proof}
(a) Between birth--death and handle-slide moments the bigon counts of $CF(\tilde\Gam,\phi_t(Z_n))$ are locally
constant; a handle-slide changes the complex by an isomorphism that is unitriangular for the action order.
(b) is Stokes' theorem for $\lambda$ along the boundary of $R_n(x)$, the two primitives agreeing at $m$ by
construction.  (c) is Proposition~\ref{prop:action}.  Numerically verified for the figure-eight, $n=1,2,3$, to
$2\cdot10^{-10}$ (\S\ref{sec:tests}).
\end{proof}
Without minimal position, (a) holds after cancelling the birth--death pairs, each of which is an explicit bigon
of explicit area.  The DLME-type level bounds follow from $|\delta_n(x)|\le\operatorname{area}$ swept; the exact
statement is (b).

\subsection{The instanton side: Floer's perturbation}
The pillowcase statement has an exact instanton counterpart, and it is Floer's original proof of the triangle
\cite{Floer,BD}.  By Sivek--Zentner \cite[Thm.~4.5]{SZ}, any area-preserving isotopy of the pillowcase applied to
the filling line is approximated by shearing isotopies induced by holonomy perturbations supported in the collar
$T^2\times I\subset Y_{1/n}$.  Choose such a perturbation $\pi_n$ moving $L_{1/n}$ onto (a smoothing of) $Z_n$.
\begin{enumerate}[label=(\roman*)]
\item The critical set of $\cs+\pi_n$ on $Y_{1/n}$ is $r^{-1}(L_{1/(n-1)})\sqcup r^{-1}(L^w_0)\sqcup\{\theta\}$, all
nondegenerate (Proposition~\ref{prop:gens} applied to $Z_n$).
\item The perturbed Chern--Simons values are the cone actions.  Precisely (Appendix~\ref{app:collar}): a shear
perturbation with Hamiltonian $H$ along $\ell_v=pa+qb$ moves the line by the time-one flow of $H\circ\ell_v$ for
$\omega_P=2\,da\wedge db$, and the perturbed functional on a perturbed-flat collar piece equals the generating function
$S=H(\ell)-\tfrac12\ell H'(\ell)$ of the shear, so that
$\cs_\pi(\rho_1)-\cs_\pi(\rho_0)\equiv\int_\gamma\lambda_0-\big[f_{L^\pi}\big]$ with $f_{L^\pi}$ the primitive of
$\lambda_0=a\,db-b\,da$ along the perturbed line: the difference is twice the signed area between $\gamma$ and $L^\pi$.
(The naive guess ``add $h(\rho_1)-h(\rho_0)$'' is wrong: the curvature in the collar contributes $-\tfrac12\ell H'(\ell)$.)
\item The perturbed differential has four blocks with respect to the two halves of the critical set.  The
\emph{literal} two-step filtration (vanishing of the wrong-way block, from $Y_0$-type to $Y_{1/(n-1)}$-type
generators) is not in Floer's or Braam--Donaldson's argument as reported by later sources, and is not implied by
energy or index (Appendix~\ref{fp:sec}, Proposition~\ref{fp:prop:crux}): it is the instanton form of the reflex
block $\beta$.  What holds (Theorem~\ref{fp:thm:cone}): unconditionally $\CI(Y_{1/n};\pi_n)\simeq\Cone(F_n)$ by
triangle detection; under (N1) (the two cobordism maps are the identity on the diagonal blocks for small neck
parameter) there is a unitriangular isomorphism of complexes onto a cone whose connecting map is the
$A\to B$ block of the perturbed differential; under (N1)--(N3) that cone is $\Cone(F_n)$.
\item The family $\cs+t\pi_n$, $t\in[0,1]$, gives a continuation map $\CI(Y_{1/n};\pi_n)\to\CI(Y_{1/n})$; when the
critical points move without birth--death (the instanton form of minimal position), it is an isomorphism of
complexes, equal to $x\mapsto x'$ up to handle-slides, and the critical values satisfy
$\cs(x')=\cs_{\pi_n}(x)+\int_0^1\big(\partial_t\cs_{t\pi_n}\big)(x_t)\,dt$, the integral being $\delta_n(x)$.
\end{enumerate}
\begin{theorem}[exact filtered cone, instanton side]\label{thm:exact-inst}
Assume \ref{H:alex}--\ref{H:Y}, transversality of $\Gam$ to $L_{1/(n-1)},L^w_0,L_{1/n}$, and minimal position.
\begin{enumerate}[label=(\alph*)]
\item \textup{(exact, unconditional)} The generators of $\CI(Y_{1/n})$ are in bijection with those of
$\Cone(F_n)$, and the Chern--Simons value of the generator corresponding to $x$ is $A_{\rm cone}(x)+\delta_n(x)+c_i$.
\item \textup{(conditional on (N1)--(N3) of Appendix~\ref{fp:sec})} $\CI(Y_{1/n})\cong\Cone(F_n)$ as
$\Z/8$-graded complexes, the isomorphism being the generator bijection composed with a unitriangular change of
basis.
\end{enumerate}
The differential of $\CI(Y_{1/n})$ is not assumed to be a bigon count.  Since the differential of $\CI(Y_{1/n})$
strictly decreases $\cs$, (a) already constrains it: with (b), the filtered complex is the cone differential
written in the basis $x'$, conjugated by a unitriangular matrix that the levels in (a) bound.
\end{theorem}
\begin{proof}[Proof sketch]
(i), (ii) are Propositions \ref{prop:gens}, \ref{prop:action} with the perturbation term (Appendix~\ref{app:collar});
(iii) is Appendix~\ref{fp:sec}; (iv) is the continuation argument with the first-variation formula; the integral
equals $\delta_n(x)$ by (ii) and Stokes.  The zero-flux condition is automatic in Floer's construction: the
perturbed line is the graph $a+nb=-g(b)$ with $g$ odd, so $\int g=0$.
\end{proof}
What remains for a complete proof of (b): conditions (N1)--(N3) of Appendix~\ref{fp:sec} (diagonal blocks of the
two cobordism maps equal to the identity for a small neck; $f_C\circ f_B=0$ on chains, not just up to homotopy;
the connecting block homotopic to $F_n$), which the sources give only up to homotopy; and the approximation of $Z_n$ by
shearing isotopies with control of the perturbed Chern--Simons values.  For the latter, Appendix~\ref{app:collar}
shows: Sivek--Zentner's approximation is $C^r$ for every fixed $r$ (not $C^\infty$), a diagonal argument (round
the corners of $Z_n$, then approximate in $C^1$) makes the actions converge, but the perturbations themselves do
not converge, so uniform compactness for the perturbed complexes along the sequence is open.  The cleanest route
avoids the limit: fix one smoothing $Z_n^\epsilon$, use the exact formula for its actions, and let $\delta_n$ absorb the
$O(\epsilon)$ difference.
